\chapter{Problem martyngałowy i twierdzenie Feynmana--Kaca}
\label{ch:martingale-fk}

\section{Od generatora do prawa procesu: motywacja}

W rozdziale \ref{ch:sde-properties} rozpoczęliśmy od procesu $X$, a następnie
wyznaczyliśmy jego generator $L$. Można odwrócić to rozumowanie:
\emph{czy operator $L$ i rozkład początkowy wyznaczają prawo procesu?}
Nie należy utożsamiać tego pytania z istnieniem silnego rozwiązania SDE.
Silne rozwiązanie wymaga zadania procesu Wienera i jego filtracji,
natomiast prawo procesu można próbować scharakteryzować bez wskazywania
konkretnego źródła szumu.

Niech $b:\R^d\to\R^d$ i $\sigma:\R^d\to\R^{d\times m}$ będą
ciągłe, a $a=\sigma\sigma^{\mathsf T}$. Rozważamy operator
\[
Lf(x)=\sum_i b_i(x)\partial_i f(x)
 +\frac12\sum_{i,j}a_{ij}(x)\partial_{ij}f(x),
 \qquad f\in C_c^2(\R^d).
\]
W równaniu SDE składnik $\int\nabla f(X_s)\sigma(X_s)\,\dd W_s$
ma średnią warunkową zero. Istotą konstrukcji jest więc fakt, że
$f(X_t)-f(X_0)-\int_0^tLf(X_s)\,\dd s$ jest martyngałem.

\begin{definicja}[Problem martyngałowy dla operatora $L$]
\label{def:mp}
Niech $\mu$ będzie miarą probabilistyczną na $\R^d$ i niech
$\Omega_C=C([0,\infty);\R^d)$ będzie przestrzenią ciągłych trajektorii
z jej $\sigma$-ciałem borelowskim (dla topologii zbieżności jednostajnej
na przedziałach zwartych). Niech $Y_t(\omega)=\omega(t)$ i
$\calF_t^Y=\sigma(Y_s:0\le s\le t)$.
Mówimy, że prawo probabilistyczne $\mathbf Q$ na $\Omega_C$
\emph{rozwiązuje problem martyngałowy $(L,\mu)$}, jeśli
$\Law_{\mathbf Q}(Y_0)=\mu$ i dla każdej $f\in C_c^2(\R^d)$
proces
\[
M_t^f=f(Y_t)-f(Y_0)-\int_0^tLf(Y_s)\,\dd s
\]
jest martyngałem względem $(\calF_t^Y)$ na każdym skończonym
przedziale czasu. Zakładamy przy tym całkowalność występujących
wyrażeń. Gdy $b,\sigma$ są lokalnie ograniczone, $Lf$ jest ograniczone
ze względu na zwarty nośnik $f$, więc całkowalność jest automatyczna.
\end{definicja}

\begin{uwaga}[Co jest, a czego nie ma w definicji?]
W definicji nie występuje proces Wienera. Warunek
martyngałowy opisuje \emph{warunkowe średnie przyrostów funkcji
stanu}. Własność ta jest silniejsza od równości
$\E f(Y_t)-\E f(Y_0)=\int_0^t\E Lf(Y_s)\,\dd s$:
ta ostatnia nie zawiera informacji o przeszłości $\calF_s^Y$.
Różne prawa mogą spełniać te same słabe równania dla rozkładów
jednowymiarowych; problem martyngałowy kontroluje całe prawo trajektorii.
\end{uwaga}

\begin{twierdzenie}[SDE rozwiązuje problem martyngałowy]
\label{th:sde-to-mp}
Niech $b,\sigma$ będą globalnie Lipschitzowskie i niech $X$ będzie
jedynym silnym rozwiązaniem
$\dd X_t=b(X_t)\,\dd t+\sigma(X_t)\,\dd W_t$
z $\Law(X_0)=\mu$ oraz $\E\norm{X_0}^2<\infty$.
Wówczas prawo procesu $X$ na $C([0,\infty);\R^d)$ rozwiązuje
problem martyngałowy $(L,\mu)$.
\end{twierdzenie}
\begin{proof}
Dla ustalonego $f\in C_c^2$ wzór Itô daje
$M_t^f=\int_0^t\nabla f(X_s)\sigma(X_s)\,\dd W_s$.
Ponieważ $\nabla f$ znika poza zwartym zbiorem, a $\sigma$ jest
tam ograniczona, integrand jest globalnie ograniczony.
Dlatego całka jest prawdziwym, kwadratowo całkowalnym
martyngałem względem filtracji napędzającej. Musimy jeszcze
uzasadnić przejście do mniejszej filtracji naturalnej.
Dla $s\le t$ i ograniczonej zmiennej $G$, mierzalnej względem
$\sigma(X_u:u\le s)$, mamy $\E[(M_t^f-M_s^f)G]=0$,
gdyż $G$ jest $\calF_s$-mierzalna. Jest to definicyjny
warunek martyngałowy względem filtracji naturalnej.
Transportując tożsamość przez odwzorowanie
$\omega\mapsto(X_t(\omega))_{t\ge0}$ otrzymujemy
wymagany warunek dla prawa kanonicznego.
\end{proof}

\section{Przykład fundamentalny: problem martyngałowy ruchu Browna}

Powyższe twierdzenie jest kierunkiem łatwiejszym: z równania
otrzymujemy problem martyngałowy. Pokażemy teraz na ważnym
przykładzie, że problem potrafi także \emph{wyznaczyć proces}.

\begin{lemat}[Rozszerzenie funkcji testowej przez obcięcie]
\label{lem:mp-cutoff}
Załóżmy, że $Y$ jest ciągłym procesem o wartościach w $\R^d$,
którego prawo rozwiązuje problem martyngałowy dla
$L=\frac12\Delta$ i rozkładu początkowego $\delta_x$.
Jeżeli $f\in C^2(\R^d)$ ma ograniczone wartości,
pochodne pierwszego i drugiego rzędu, to
$f(Y_t)-f(Y_0)-\frac12\int_0^t\Delta f(Y_s)\,\dd s$
jest martyngałem na każdym przedziale $[0,T]$.
\end{lemat}
\begin{proof}
Wybierzmy $\chi_R\in C_c^\infty(\R^d)$ równą $1$ na
kuli promienia $R$ i $0$ poza kulą promienia $2R$.
Funkcja $f_R=\chi_R f$ jest klasy $C_c^2$.
Niech $\tau_R=\inf\{t\ge0:\norm{Y_t}\ge R\}$.
Dla $t\le\tau_R$ zarówno $f_R(Y_t)=f(Y_t)$,
jak i $\Delta f_R(Y_t)=\Delta f(Y_t)$.
Stąd proces zatrzymany w $\tau_R$ jest martyngałem,
przez opcjonalne zatrzymywanie w ograniczonym czasie
$T\wedge\tau_R$. Dla każdej ciągłej trajektorii
$\sup_{0\le t\le T}\norm{Y_t}<\infty$, więc
$\tau_R>T$ od pewnego $R$ (zależnego od trajektorii).
Ponadto na $[0,T]$ zmienne
$|f(Y_{t\wedge\tau_R})-f(Y_0)
-\frac12\int_0^{t\wedge\tau_R}\Delta f(Y_s)\,\dd s|$
są ograniczone przez
$2\norm f_\infty+\frac T2\norm{\Delta f}_\infty$.
Dla każdego $A\in\calF_s^Y$ przechodzimy więc
do granicy w tożsamości
$\E[\ind A(M_{t\wedge\tau_R}^f-M_{s\wedge\tau_R}^f)]=0$
za pomocą zbieżności zdominowanej. Otrzymujemy tezę.
\end{proof}

\begin{twierdzenie}[Jednoznaczność problemu martyngałowego dla ruchu Browna]
\label{th:mp-bm}
Niech $x\in\R^d$ i $L=\frac12\Delta$ na $C_c^2(\R^d)$.
Jedynym prawem na $C([0,\infty);\R^d)$ rozwiązującym
problem martyngałowy $(L,\delta_x)$ jest prawo procesu
$x+W$, gdzie $W$ jest standardowym $d$-wymiarowym
ruchem Browna.
\end{twierdzenie}
\begin{proof}
Istnienie wynika z twierdzenia \ref{th:sde-to-mp},
przy $b=0$, $\sigma=I_d$. Dla jednoznaczności niech
$Y$ będzie procesem kanonicznym pod dowolnym innym
rozwiązaniem $\mathbf Q$. Ustalmy $u\in\R^d$ i
$f_u(y)=e^{i\langle u,y\rangle}$. W rzeczywistym
problemie martyngałowym stosujemy osobno części rzeczywistą
i urojoną. Pochodne $f_u$ są ograniczone, zatem lemat
\ref{lem:mp-cutoff} daje, że
\[
Z_t=f_u(Y_t)+\frac{\norm u^2}{2}\int_0^t f_u(Y_s)\,\dd s
\]
jest martyngałem zespolonym. Zróżniczkowanie iloczynu
$e^{\norm u^2t/2}f_u(Y_t)$ z deterministyczną funkcją
czasu (całkowanie przez części) pokazuje, że
$H_t=e^{\norm u^2t/2}f_u(Y_t)$ jest martyngałem.
Na $[0,T]$ jest on ograniczony przez $e^{\norm u^2T/2}$,
nie ma więc problemu z lokalizacją.
Dla $s\le t$ otrzymujemy
\[
\E_{\mathbf Q}[e^{i\langle u,Y_t-Y_s\rangle}\mid\calF_s^Y]
=e^{-\norm u^2(t-s)/2}.
\]
Lewa strona jest warunkową funkcją charakterystyczną
przyrostu. Prawo przyrostu jest więc normalne
$\Normal(0,(t-s)I_d)$ i nie zależy od $\calF_s^Y$.
Jednoznaczność funkcji charakterystycznych i
argument z klasą monotoniczną dowodzą niezależności
przyrostu od całej przeszłości. Ponieważ $Y_0=x$ i
trajektorie są ciągłe z definicji przestrzeni kanonicznej,
$Y-x$ jest standardowym ruchem Browna.
Prawo takiego procesu jest wyznaczone przez rozkłady
skończeniewymiarowe, co kończy dowód.
\end{proof}

\begin{przyklad}[Dlaczego testujemy wiele funkcji?]
Sam warunek $\E(Y_t-Y_s\mid\calF_s)=0$ nie wystarcza
do rozpoznania procesu Wienera: spełnia go również
wiele procesów skokowych i ciągłych martyngałów o innej
wariacji kwadratowej. Funkcje
$e^{i\langle u,x\rangle}$ kodują wszystkie
funkcje charakterystyczne przyrostów i pozwalają
rozpoznać całe ich rozkłady.
\end{przyklad}

\section{Półgrupy, potencjały i motywacja formuły Feynmana--Kaca}

Niech $X^x$ rozwiązuje jednorodne SDE o globalnie
Lipschitzowskich współczynnikach, jak w
rozdziale \ref{ch:sde-properties}. Dla ograniczonej
funkcji borelowskiej $f$ i $t\ge0$ znamy już półgrupę
$P_tf(x)=\E_xf(X_t)$. Jeżeli cząstka może zostać
\emph{zabita} z intensywnością zależną od położenia,
to każdą trajektorię trzeba dodatkowo ważyć przez
$e^{-\int_0^tV(X_s)\,\dd s}$. Prowadzi to do
nowej rodziny operatorów.

\begin{definicja}[Półgrupa Feynmana--Kaca]
\label{def:fk-semigroup}
Niech $V:\R^d\to\R$ będzie ograniczoną funkcją
borelowską. Dla ograniczonej borelowskiej $f$ określamy
\[
(T_t^Vf)(x)=\E_x\left[
\exp\left(-\int_0^tV(X_s)\,\dd s\right)f(X_t)\right].
\]
Całka w wykładniku jest zwykłą całką względem czasu.
Cały nawias oznacza funkcjonał losowej trajektorii.
\end{definicja}

\begin{twierdzenie}[Własność półgrupowa i oszacowanie]
\label{th:fk-semigroup}
Przy powyższych założeniach $T_0^V=I$,
$T_{s+t}^V=T_s^VT_t^V$ na ograniczonych
funkcjach borelowskich oraz
$\norm{T_t^Vf}_\infty\le e^{t\norm V_\infty}\norm f_\infty$.
Jeśli $V\ge0$, to $T_t^V$ jest dodatnią kontrakcją:
$0\le f\le1$ implikuje $0\le T_t^Vf\le1$.
\end{twierdzenie}
\begin{proof}
Oszacowanie wynika z nierówności
$|e^{-\int_0^t V(X_s)\,\dd s}|\le e^{t\norm V_\infty}$.
Dla własności półgrupowej rozbijamy
$\int_0^{s+t}V(X_r)\,\dd r$ na odcinki
$[0,s]$ i $[s,s+t]$. Warunkowo względem $\calF_s$
dalszy fragment procesu, po ustaleniu $X_s$, ma
prawo procesu startującego w $X_s$.
Stosujemy więc własność Markowa nie tylko do
jednego punktu końcowego, lecz do ograniczonego
mierzalnego funkcjonału całego fragmentu przyszłej trajektorii.
Wyjaśnijmy przejście, które jest istotne przy borelowskim $V$.
Dla każdej skończonej rodziny chwil $r_1,\ldots,r_k\in[0,t]$
własność Markowa i iterowanie wartości oczekiwanej warunkowej
dają stosowną tożsamość dla funkcji
$F(w)=\varphi(w(r_1),\ldots,w(r_k))$.
Zbiory takich funkcji cylindrycznych generują borelowskie
$\sigma$-ciało na $C([0,t];\R^d)$; twierdzenie o klasie
monotonicznej rozszerza tożsamość na wszystkie ograniczone
borelowskie funkcjonały $F$. Funkcjonał
$F(w)=\exp(-\int_0^tV(w(r))\,\dd r)f(w(t))$
jest borelowski: odwzorowanie $(r,w)\mapsto V(w(r))$
jest wspólnie mierzalne, a całkowanie ograniczonej funkcji
mierzalnej względem $r$ zachowuje mierzalność względem $w$
(z twierdzenia Fubiniego). Nie trzeba więc zakładać
zbieżności sum Riemanna dla potencjału jedynie borelowskiego.
Dostajemy
\[
\E_x\left[e^{-\int_0^{s+t}V(X_r)\,\dd r}f(X_{s+t})
\mid\calF_s\right]
=e^{-\int_0^sV(X_r)\,\dd r}(T_t^Vf)(X_s).
\]
Po wzięciu wartości oczekiwanych dostajemy
$T_{s+t}^Vf=T_s^V(T_t^Vf)$. Dodatniość dla $V\ge0$
jest bezpośrednia.
\end{proof}

\section{Wzór Feynmana--Kaca dla klasycznych rozwiązań PDE}

Poniższe twierdzenie ma celowo precyzyjny zakres.
Najpierw \emph{weryfikujemy} reprezentację dla
istniejącego klasycznego rozwiązania równania wstecznego.
Samo istnienie rozwiązania $C^{1,2}$ dla ogólnych
współczynników wymaga odrębnych twierdzeń teorii PDE.

\begin{twierdzenie}[Feynman--Kac: postać wsteczna z wymuszeniem]
\label{th:feynman-kac}
Niech $T>0$, $b:\R^d\to\R^d$,
$\sigma:\R^d\to\R^{d\times m}$ będą globalnie
Lipschitzowskie, $V:\R^d\to\R$ oraz
$q:[0,T]\times\R^d\to\R$ będą ograniczonymi funkcjami
ciągłymi, a $g:\R^d\to\R$ będzie ograniczona
i ciągła. Niech $u\in C^{1,2}([0,T)\times\R^d)
\cap C([0,T]\times\R^d)$ będzie ograniczona,
a jej pochodne pierwszego i drugiego rzędu względem
$x$ oraz pierwsza pochodna względem $t$ będą
ograniczone na $[0,T)\times\R^d$.
Załóżmy, że $u$ spełnia
\[
\partial_tu(t,x)+Lu(t,x)-V(x)u(t,x)+q(t,x)=0
\quad (t<T),\qquad u(T,x)=g(x).
\]
Wówczas dla każdego $t\in[0,T]$ i $x\in\R^d$
\[
\begin{split}
u(t,x)=\E\Big[&e^{-\int_t^T V(X_r^{t,x})\,\dd r}
                         g(X_T^{t,x})\\
 &+\int_t^T e^{-\int_t^s V(X_r^{t,x})\,\dd r}
                         q(s,X_s^{t,x})\,\dd s\Big],
\end{split}
\tag{7.1}\label{eq:fk}
\]
gdzie $(X_s^{t,x})_{s\ge t}$ jest rozwiązaniem SDE
startującym w $x$ w chwili $t$.
W klasie funkcji o podanej regularności i ograniczoności
rozwiązanie równania wstecznego jest jednoznaczne.
\end{twierdzenie}
\begin{proof}
Niech $X_s=X_s^{t,x}$ oraz
$A_s=\exp(-\int_t^sV(X_r)\,\dd r)$ dla $s\in[t,T]$.
Jest to ciągły proces o skończonej wariacji, spełniający
$\dd A_s=-V(X_s)A_s\,\dd s$, $A_t=1$.
Dla $s<T$ stosujemy wzór Itô do $u(s,X_s)$
i regułę iloczynu dla $A_su(s,X_s)$.
Kowariacja $A$ z ciągłym lokalnym martyngałem wynosi zero,
więc
\[
\dd[A_su(s,X_s)]
=A_s(\partial_su+Lu-Vu)(s,X_s)\,\dd s
 +A_s\nabla u(s,X_s)\sigma(X_s)\,\dd W_s.
\]
Po podstawieniu PDE dryf wynosi $-A_sq(s,X_s)\,\dd s$.
Dlatego proces
$A_su(s,X_s)+\int_t^sA_rq(r,X_r)\,\dd r$
jest lokalnym martyngałem. W rzeczywistości jest
prawdziwym martyngałem: wszystkie jego składniki poza
całką stochastyczną są ograniczone przez
$e^{T\norm V_\infty}(\norm u_\infty+T\norm q_\infty)$,
a sama całka ma z równania tę samą całkowalną,
a nawet ograniczoną na $[t,T)$ różnicę tych składników.
Po zatrzymaniu i zastosowaniu zbieżności zdominowanej
usuwamy lokalizację. Biorąc wartości oczekiwane
w chwili $T-\varepsilon$ i przechodząc $\varepsilon\downarrow0$
przez ciągłość $u$, $X$ i ograniczoność $A,q$, dostajemy
\eqref{eq:fk}.
Jeśli $u_1,u_2$ spełniają te same warunki, to oba
mają tę samą reprezentację prawostronną, a więc są równe.
\end{proof}

\begin{uwaga}[Wersja z potencjałem zależnym od czasu]
Ten sam dowód działa dla ograniczonej ciągłej funkcji
$V:[0,T]\times\R^d\to\R$, po zamianie
$V(X_s)$ na $V(s,X_s)$ w całce oraz w PDE.
Nie trzeba przy tym zakładać, że $V$ jest nieujemna;
ograniczoność zapewnia całkowalność czynnika wykładniczego.
\end{uwaga}

\section{Przykłady rachunkowe i rozwiązane zadania}

\begin{zadanie}[Równanie ciepła i jądro gaussowskie]
Niech $g\in C_b^2(\R)$, przy czym $g''$ jest
jednostajnie ciągła. Rozwiązać problem
$\partial_tu+\frac12\partial_{xx}u=0$ dla $t<T$
i $u(T,x)=g(x)$.

\emph{Rozwiązanie.} Weźmy $X_s^{t,x}=x+W_s-W_t$.
Kandydatem jest
\[
u(t,x)=\E g(x+W_{T-t})
=\int_\R g(y)\frac{e^{-(y-x)^2/(2(T-t))}}
{\sqrt{2\pi(T-t)}}\,\dd\lambda_1(y)
\quad(t<T).
\]
Dla $h>0$ można różniczkować pod całką względem $x$,
gdyż $g',g''$ są ograniczone: $u_x=\E g'(x+W_h)$,
$u_{xx}=\E g''(x+W_h)$.
Z własności półgrupy ciepła oraz jednostajnej ciągłości
$g''$ wynika
$\partial_h\E g(x+W_h)=\frac12\E g''(x+W_h)$.
Można to też otrzymać z wzoru Itô i zbieżności
zdominowanej po podzieleniu przez mały przyrost czasu.
Ponieważ $h=T-t$, mamy $u_t=-\frac12u_{xx}$.
Jednostajna ciągłość $g$ zapewnia warunek końcowy.
Zatem $u$ jest klasycznym rozwiązaniem z twierdzenia
\ref{th:feynman-kac}; otrzymaliśmy zarówno istnienie,
jak i jednoznaczność w rozważanej klasie.
\end{zadanie}

\begin{zadanie}[Stały potencjał: zabijanie trajektorii]
Niech $\kappa\ge0$, $V(x)=\kappa$, $q=0$ i $g\equiv1$.
Wykazać, że rozwiązaniem równania
$\partial_tu+\frac12\Delta u-\kappa u=0$,
$u(T,x)=1$, jest $u(t,x)=e^{-\kappa(T-t)}$.

\emph{Rozwiązanie.} Dla każdej trajektorii
$\int_t^TV(X_s)\,\dd s=\kappa(T-t)$.
Wzór \eqref{eq:fk} daje więc
$u(t,x)=\E e^{-\kappa(T-t)}=e^{-\kappa(T-t)}$.
Weryfikacja PDE jest bezpośrednia:
$u_t=\kappa u$, $\Delta u=0$, stąd
$u_t-\kappa u=0$. W interpretacji probabilistycznej
$u$ jest prawdopodobieństwem przeżycia niezależnego
czasu wykładniczego o intensywności $\kappa$ od $t$ do $T$.
\end{zadanie}

\begin{zadanie}[Równanie Poissona i średni czas wyjścia]
Niech $B_t$ będzie standardowym ruchem Browna w $\R^d$,
a $\tau_R=\inf\{t\ge0:\norm{x+B_t}\ge R\}$ dla $\norm x<R$.
Wyznaczyć $\E_x\tau_R$ i wyjaśnić, czemu równanie
$Lu=-1$ z warunkiem $u=0$ na sferze daje tę wielkość.

\emph{Rozwiązanie.} Funkcja
$u(x)=(R^2-\norm x^2)/d$ spełnia
$\frac12\Delta u=-1$ w kuli oraz $u=0$ na jej brzegu.
Z wzoru Itô dla procesu zatrzymanego wynika
$u(x+B_{t\wedge\tau_R})=u(x)-(t\wedge\tau_R)+M_t$,
gdzie $M_t$ jest całką stochastyczną z ograniczonym
integrandem na zatrzymanym procesie (bo $\nabla u$ jest
ograniczony w domkniętej kuli). Stąd
$\E(t\wedge\tau_R)\le u(x)$,
ponieważ $u\ge0$ na kuli. Zatem $\E\tau_R<\infty$.
Po przejściu $t\to\infty$ przez zbieżność zdominowaną
w pierwszym składniku oraz monotoniczną w drugim
otrzymujemy $u(x)=\E_x\tau_R$.
To ten sam mechanizm co we wzorze Feynmana--Kaca:
obserwowaną wielkość przedstawiamy przez funkcjonał
procesu i identyfikujemy odpowiedni problem brzegowy.
\end{zadanie}

\begin{center}
\begin{tikzpicture}[>=Latex,node distance=7mm]
\node[draw,rounded corners,fill=blue!6,align=center,minimum width=31mm] (sde)
{równanie SDE\\proces $X$};
\node[draw,rounded corners,fill=orange!8,align=center,right=13mm of sde,minimum width=34mm] (gen)
{generator $L$\\problem martyngałowy};
\node[draw,rounded corners,fill=green!7,align=center,right=13mm of gen,minimum width=31mm] (pde)
{równanie PDE\\wzór Feynmana--Kaca};
\draw[->,thick] (sde)--(gen);
\draw[->,thick] (gen)--(pde);
\draw[->,bend left=35,thin,dashed] (pde.north) to node[above,scale=.78]
{reprezentacja rozwiązania} (sde.north);
\end{tikzpicture}
\par\smallskip\small Związki między procesem, jego generatorem i równaniem wstecznym.
\end{center}

\begin{uwaga}[Prawidłowy zakres wyniku o rozwiązaniach PDE]
Twierdzenie \ref{th:feynman-kac} ma charakter weryfikacyjny:
każde rozwiązanie w wymienionej klasie ma podaną reprezentację.
Sama mierzalność lub ograniczoność potencjału nie zapewnia
istnienia klasycznego rozwiązania $C^{1,2}$ bez dodatkowej
teorii regularności PDE. Podobnie równanie słabe dla rozkładów
jednowymiarowych nie zastępuje dobrze postawionego problemu
martyngałowego dla prawa całego procesu.
\end{uwaga}

\section{Granice uzyskanej teorii}
Wykazaliśmy, że silne rozwiązanie SDE spełnia problem martyngałowy,
a dla operatora $\frac12\Delta$ udowodniliśmy także kierunek odwrotny
na poziomie prawa procesu. W pełnej teorii związek między dobrze
postawionym problemem martyngałowym a słabymi rozwiązaniami SDE
wymaga dalszych narzędzi, m.in. konstrukcji całek względem ogólnych
ciągłych lokalnych martyngałów. Rozwiniemy je po wprowadzeniu
półmartyngałów. Wzór Feynmana--Kaca został udowodniony jako
twierdzenie weryfikacyjne dla klasycznych rozwiązań PDE;
nie przypisujemy mu bez dodatkowych twierdzeń regularności,
której nie udowodniliśmy.


% ============================================================================
% ROZDZIAL 8 — Zmiana miary i Girsanov
% ============================================================================
